\scanpage{25}
\mspage{\handwritten{16}}
définissant des hypersurfaces $H_o,..H_{n-1}$ telles que $X\handwritten{\cap}H_o\handwritten{\cap}..\handwritten{\cap}H_{n-1}=Z$ soit un schéma fini \inserted{\handwritten{étale}} sur $k$, de degré $rd_1^n$, dont tous\deleted{\unsure{es}} les points sont rationnels sur $k$\inserted{et distinct des $x_i$.}. Quitte à augmenter $k$ et $d_1$, on peut supposer qu'il existe une forme $\handwritten{\phi}_n$ de même degré $d_1$, définissant $H_n$, tel\handwritten{\unsure{le|l|e}} que $H_n\handwritten{\cap}Z$ soit de cardinal $N$. Pour tout entier $d_2\handwritten{\geq}n+1$, on sait qu'il existe une forme \handwritten{$\underline{\phi}$} à $n+1$ variables \deleted{telle que} définie sur $k$, de degré $d_2$, \deleted{telle que} n'ayant pas de zéro non trivial sur $k$. Alors
\[
\handwritten{\phi}=\handwritten{\underline{\phi}}\,(\handwritten{\phi}_o,\ldots,\handwritten{\phi}_n)
\]
est une forme de degré $d_2d_1$ sur $P^r_k$, telle que $(H_{\handwritten{\phi}}\handwritten{\cap}X)(k)=(H_n\handwritten{\cap}Z)(k)$ soit de cardinal $N$. De plus, comme cet ensemble ne contient pas les $x_i$, \handwritten{$\phi$} définit bien un point de $U_{d_1d_2}$. Lorsque $N$ est 0 ou 1\inserted{(plus généralement $N\handwritten{\leq}r$)}, on a terminé, car on peut prendre $d_1=1$, et 7.6 marche avec $d'_o=n+1$. Si on a $N>r$, on note que pour tout $d_{\handwritten{3}}\handwritten{\geq}n+1$, il existe une forme \handwritten{$\unsure{\psi|\mu}$} de degré $d_{\handwritten{3}}$ \inserted{rat/$k$} telle que $(H_{\handwritten{\psi}}\handwritten{\cap}X)(k)=\handwritten{\varnothing}$ d'après ce qu'on vient de voir, et alors la forme \handwritten{$\phi\ \psi$} est de degré \inserted{$d=$} $d_1d_2+d_{\handwritten{3}}$ et $(H_{\handwritten{\phi\psi}}\handwritten{\cap}X)(k)=Z(k)$ a comme cardinal $N$, et \handwritten{$\phi\psi$} définit encore un élément de $U_d$. On peut donc prendre alors $d'_o=d_1\handwritten{\unsure{d|\phi}}_2+(n+1)=(d_1+1)(n+1)$ (en prenant, comme il est loisible, $d_2=n+1$.

\uline{Remarques} 7.7. \handwritten{a)} On trouve ainsi une estimation assez précise de $d'_o$ de 7.6 \deleted{\unsure{(donc}} en fonction de $N$ et de $r$, donc du $d_o$ de 7.1 en fonction de $N=p-1$ et de $r$, \deleted{\unsure{lorsqu'on désire dans}} et du plus petit entier rendant valables les égalités (*). Lorsque dans la conclusion de 7.1 on se contente d'affirmer que les éléments de $G_o$ donnent lieu à au moins deux traces distinctes (ce qui suffit à conclure que $G_o$ n'est pas un $p$-groupe et a fortiori n'est pas le groupe unité, donc $\dim V\handwritten{\geq}1$ i.e. frobenius sur $V$ n'est pas nilpotent) on peut prendre \deleted{pour} simplement $d'_o=n+1$, et pour $d_o$ le plus petit entier $\handwritten{\geq}n+1$ tel que l'on ait (*). Ce choix de $d_o$ est suffisant aussi, par conséquent, pour avoir la validité des
% archive sheet number: 24
